\section*{Hoofdstuk \ref{ch:g6:fractions} --- Breuken: eerste stappen}

\begin{solution}{exo:g6:fractions:1}
Cake: $\frac38$ opgegeten. \quad Chocolade: $\frac{7}{12}$. \quad
Taart: $\frac66 = 1$ (de hele taart).
\end{solution}

\begin{solution}{exo:g6:fractions:2}
Vier van de vijf delen zijn gekleurd. $\frac45 < 1$ (vier van de vijf delen is
minder dan de hele strook); $\frac65 > 1$ (zes vijfden is één hele strook en
een vijfde erbij).
\end{solution}

\begin{solution}{exo:g6:fractions:3}
Op de lijn met streepjes per derde: $\frac13$ ligt één streepje na $0$;
$\frac53$ ligt één streepje voor $2$ (want $2 = \frac63$); $\frac73$ ligt één
streepje na $2$. Dus liggen $\frac53$ en $\frac73$ beide op afstand $\frac13$
van $2$.
\end{solution}

\begin{solution}{exo:g6:fractions:4}
$\frac94 = 2 + \frac14$; \quad
$\frac{13}{5} = 2 + \frac35$; \quad
$\frac{12}{6} = 2$ (precies); \quad
$\frac{25}{8} = 3 + \frac18$.
\end{solution}

\begin{solution}{exo:g6:fractions:5}
$\frac23 = \frac{8}{12}$ (met $4$ vermenigvuldigd); \quad
$\frac{15}{20} = \frac34$ (door $5$ gedeeld); \quad
$\frac{7}{10} = \frac{21}{30}$ (met $3$ vermenigvuldigd); \quad
$\frac56 = \frac{10}{12}$ (door $2$ gedeeld).
\end{solution}

\begin{solution}{exo:g6:fractions:6}
$\frac68 = \frac34$; \quad
$\frac{15}{25} = \frac35$; \quad
$\frac{18}{24} = \frac{9}{12} = \frac34$; \quad
$\frac{35}{7} = 5$ (een heel getal).
\end{solution}

\begin{solution}{exo:g6:fractions:7}
$\frac12$ van $86$: $86 \div 2 = 43$.

$\frac34$ van $60$: $60 \div 4 = 15$, dan $15 \times 3 = 45$.

$\frac25$ van $35$: $35 \div 5 = 7$, dan $7 \times 2 = 14$.

$\frac{7}{10}$ van $250$: $250 \div 10 = 25$, dan $25 \times 7 = 175$.
\end{solution}

\begin{solution}{exo:g6:fractions:8}
$\frac12 > \frac13$: \omterm{def:g3:division:half}{helften} zijn grotere delen dan derden.
$\frac35 > \frac25$: dezelfde delen, en meer ervan.
$\frac34 = \frac68 < \frac78$: zeven achtsten verslaan zes achtsten.
\end{solution}

\begin{solution}{exo:g6:fractions:9}
$\frac12 = 0.5$; \quad $\frac34 = 0.75$; \quad $\frac{7}{10} = 0.7$; \quad
$\frac94 = 2.25$. De uitzondering is $\frac13 = 0.333\dots$, waarvan de
decimale schrijfwijze nooit ophoudt: dat is geen \omterm{def:g6:decimals:places}{decimaal getal}.
\end{solution}

\begin{solution}{exo:g6:fractions:10}
Met de bus: $\frac34$ van $28$: $28 \div 4 = 7$, dan $7 \times 3 = 21$
leerlingen. Lopend: $28 - 21 = 7$ leerlingen (dat is het overblijvende \omterm{def:g3:division:half}{kwart}:
$28 \div 4 = 7$). Controle: $21 + 7 = 28$.
\end{solution}

\begin{solution}{exo:g6:fractions:11}
Twee derde van het spaargeld is $18$ euro, dus is \emph{één} derde
$18 \div 2 = 9$ euro, en was het hele spaargeld (drie derden)
$9 \times 3 = 27$ euro. Controle: $\frac23$ van $27$ is $18$.
\end{solution}

\begin{solution}{exo:g6:fractions:12}
Het water ging van $\frac58$ naar $\frac48$ (een \omterm{def:g3:division:half}{helft} is vier achtsten) van de
tank: die $10$ liter is dus $\frac58 - \frac48 = \frac18$ van de \omterm{def:g2:measure:units}{inhoud}. De
volle tank bevat dus $10 \times 8 = 80$ liter. Controle: $\frac58$ van $80$ is
$50$, min $10$ blijft $40$, en dat is inderdaad de \omterm{def:g3:division:half}{helft} van $80$.
\end{solution}

\begin{solution}{pb:g6:fractions:1}
\textbf{1.} De \omterm{def:g3:division:half}{helft} van $16$ blokjes is $8$ blokjes. Als \omterm{def:g6:fractions:def}{breuk}: $\frac12$ van
de reep, en omdat elke \omterm{def:g3:division:half}{helft} $8$ van de $16$ blokjes is, ook $\frac{8}{16}$ ---
dezelfde \omterm{def:g6:fractions:def}{breuk}, op twee manieren geschreven
(\cref{prop:g6:fractions:equal}).

\textbf{2.} De \omterm{def:g3:division:half}{helft} van de $8$ overgebleven blokjes is $4$ blokjes, en dat is
$\frac{4}{16}$ van de reep. Vereenvoudigd door $4$: $\frac{4}{16} = \frac14$.
Op het plaatje: de overgebleven \omterm{def:g3:division:half}{helft} in twee knippen geeft twee \omterm{def:g3:division:half}{kwarten} van de
oorspronkelijke reep --- de \omterm{def:g3:division:half}{helft} van een \omterm{def:g3:division:half}{helft} is een \omterm{def:g3:division:half}{kwart}.

\textbf{3.} Derde hap: de \omterm{def:g3:division:half}{helft} van $4$ blokjes is $2$ blokjes, dus
$\frac{2}{16} = \frac18$ van de reep. Vierde hap: $1$ blokje, dus
$\frac{1}{16}$. Van de ene hap naar de volgende \emph{verdubbelt} de \omterm{def:g4:fractions:def}{noemer}:
$2$, $4$, $8$, $16$.

\textbf{4.} Opgegeten: $8 + 4 + 2 + 1 = 15$ van de $16$ blokjes, dus
$\frac{15}{16}$ van de reep. Over: $1$ blokje, $\frac{1}{16}$.

\textbf{5.} $\frac12 = \frac{8}{16}$ (\omterm{def:g4:fractions:def}{teller} en \omterm{def:g4:fractions:def}{noemer} met $8$
vermenigvuldigd), $\frac14 = \frac{4}{16}$ (met $4$),
$\frac18 = \frac{2}{16}$ (met $2$), en $\frac{1}{16}$ blijft staan. In
zestienden geteld: $8 + 4 + 2 + 1 = 15$ ervan, dus is het totaal
$\frac{15}{16}$.

\textbf{6.} Happen: $32$, $16$, $8$, $4$, $2$, $1$ blokjes. Totaal opgegeten:
$32 + 16 + 8 + 4 + 2 + 1 = 63$ blokjes, dus blijft er $1$ blokje over:
$\frac{1}{64}$ van de reep.

\textbf{7.} Na elke hap blijft er
\[
\frac12,\quad \frac14,\quad \frac18,\quad \frac{1}{16},\quad
\frac{1}{32},\quad \frac{1}{64}
\]
over: de \omterm{def:g3:division:remainder}{rest} halveert elke keer --- haar \omterm{def:g4:fractions:def}{noemer} verdubbelt. Een zevende,
verbeelde hap zou $\frac{1}{128}$ overlaten.

\textbf{8.} Elke hap eet maar \emph{de \omterm{def:g3:division:half}{helft}} van wat er over is, dus staat de
andere \omterm{def:g3:division:half}{helft} van de \omterm{def:g3:division:remainder}{rest} er nog: na elk aantal happen blijft er altijd iets
over. Het opgegeten totaal blijft dus altijd onder $1$ --- precies zoals de
twintig negens van \cref{pb:g6:decimals:1}, die steeds dichter bij $3$ komen
zonder het te bereiken.

\textbf{9.} Beide \omterm{def:g6:fractions:def}{breuken} hebben \omterm{def:g4:fractions:def}{teller} $1$, en in $128$ delen knippen geeft
kleinere delen dan in $100$ delen knippen: dus $\frac{1}{128} < \frac{1}{100}$.
Na de zevende hap is de \omterm{def:g3:division:remainder}{rest} $\frac{1}{128}$, en dus is het opgegeten deel meer
dan $1 - \frac{1}{100} = \frac{99}{100}$: bij hap zeven is de uitdaging
gewonnen.

\textbf{10.} In zestienden zijn de totalen $\frac{8}{16}$, $\frac{12}{16}$,
$\frac{14}{16}$, $\frac{15}{16}$. Elk nieuw totaal landt precies
\emph{halverwege} tussen het vorige totaal en $1$: de trap halveert steeds de
afstand die nog naar $1$ overblijft, zonder er ooit op te stappen.

\textbf{11.} Wat juist is: in elke fase van Zeno's beschrijving blijft er
inderdaad een stukje afstand over (vraag 8) --- de lijst fasen houdt nooit op.
De valstrik: de fasen worden extreem kort, in afstand \emph{en} in looptijd; de
hardloper doet niet over elke fase even lang. Dat je de loop in oneindig veel
krimpende stukjes beschrijft, maakt de loop zelf niet eindeloos: de hardloper
bereikt de muur, en vraag 9 laat zien dat de beschrijving elk punt vóór de muur
voorbijgaat.

\textbf{12.} Snijd $8$ reepjes in \omterm{def:g3:division:half}{helften}: $16$ halve stukken. Snijd $4$
reepjes in \omterm{def:g3:division:half}{kwarten}: $16$ kwartstukken. Snijd $2$ reepjes in achtsten: $16$
achtste stukken. Snijd het laatste reepje in zestienden: $16$ kleine stukjes.
Dat gebruikt $8 + 4 + 2 + 1 = 15$ reepjes, en elk kind krijgt één stuk van elke
maat.

\textbf{13.} Het deel van elk kind is
\[
\frac12 + \frac14 + \frac18 + \frac{1}{16} = \frac{15}{16}
\]
van een reep (vraag 5). Zestien zulke delen maken
$16 \times \frac{15}{16} = 15$ reepjes
(\cref{thm:g6:fractions:quotient}): precies wat er gesneden is, zonder \omterm{def:g3:division:remainder}{rest} ---
een eerlijke verdeling van $15$ reepjes over $16$ kinderen, met alleen
halveringen.

\textbf{14.} De buren krijgen elk $\frac78$ reep. Met \omterm{def:g4:fractions:def}{noemer} $16$:
$\frac78 = \frac{14}{16}$, terwijl onze kinderen $\frac{15}{16}$ krijgen. Een
kind van de $16$ krijgt dus meer: $\frac{15}{16} > \frac{14}{16}$, en wel één
zestiende reep.

\textbf{15.} In het vierkantsplaatje vult elke gekleurde hap de \omterm{def:g3:division:half}{helft} van het
nog ongekleurde hoekje, en blijft dat hoekje krimpen: na elke hap is het half
zo groot als daarvoor, en het verdwijnt nooit. De twee ware feiten: (i) na elk
eindig aantal happen blijft er een \omterm{def:g3:division:remainder}{rest} --- het opgegeten totaal is altijd
strikt kleiner dan $1$; (ii) die \omterm{def:g3:division:remainder}{rest} wordt kleiner dan elke \omterm{def:g6:fractions:def}{breuk} die je maar
wilt noemen ($\frac{1}{100}$, $\frac{1}{1000}$, \dots), als je lang genoeg
hapt. De woorden ``de oneindige \omterm{def:g1:addition:def}{som} is gelijk aan $1$'' hun precieze betekenis
geven, is het werk van de \emph{limieten}, in het bovenbouwvolume.
\end{solution}
